%% ============================================================
%%  PRAKTIKA-TXOSTENA · Robotika eta Kontrol Adimendunak
%%  Informatika Fakultateko GrAL txantiloian oinarritua (bertsio arina).
%%  Konpilatzeko: pdfLaTeX + BibTeX (Overleaf-en lehenetsitakoa).
%%
%%  \bete{...}       -> zuek idatzi beharrekoa (grisez agertzen da)
%%  \egiaztatu{...}  -> zuen kodearekin edo datuekin egiaztatu beharreko balioa (laranjaz)
%%  Entregatu aurretik, bilatu "bete{" eta "egiaztatu{" (Ctrl+F): ez du bat ere geratu behar.
%% ============================================================
\documentclass[a4paper, 11pt, oneside]{memoir}
\semiisopage
\chapterstyle{veelo}
% Euskaraz "3. kapitulua" esaten da, ez "Kapitulua 3": zenbakia bakarrik erakusten da.
\renewcommand*{\printchaptername}{\chapnamefont}

%% ---------- AZALEKO DATUAK (bete) ----------
\newcommand{\egileak}{Izen Abizen \\ Izen Abizen}
\newcommand{\ikasketak}{Informatika Ingeniaritzako Gradua}
\newcommand{\irakasgaia}{Robotika eta Kontrol Adimendunak}
\newcommand{\irakaslea}{Elena Lazkano}
\newcommand{\izenburua}{Oztopoak ekiditen eta errepideari jarraitzen}
\newcommand{\azpititulua}{Sentsoreen funtzionamendua lantzen: ROSMASTER X3 robotarekin egindako praktika}
\newcommand{\data}{\today}

\input config/makroak

\title{\izenburua}
\author{\egileak}
\date{\data}

\begin{document}

%% ---------- AZALA ----------
\input config/azala
\clearpage

%% ---------- HASIERAKO ORRIAK ----------
\frontmatter
\aliaspagestyle{chapter}{plain}
\aliaspagestyle{cleared}{plain}
\pagestyle{plain}

\chapter*{Laburpena}
\input kapituluak/laburpena
\clearpage

\pdfbookmark[1]{\contentsname}{toc}
\begin{KeepFromToc}
  \tableofcontents
\end{KeepFromToc}
\clearpage
\listoffigures
\clearpage
\listoftables

%% ---------- EDUKIA ----------
\clearpage
\mainmatter
\pagestyle{ruled}

\chapter{Sarrera}\label{ch:sarrera}
\input kapituluak/01_sarrera

\chapter{Oztopoak ekiditen}\label{ch:oztopoak}
\input kapituluak/02_oztopoak

\chapter{Errepideari jarraitzen}\label{ch:errepidea}
\input kapituluak/03_errepidea

\chapter{Ondorioak}\label{ch:ondorioak}
\input kapituluak/04_ondorioak

%% ---------- BIBLIOGRAFIA ----------
\bibliographystyle{unsrt}
\bibliography{erreferentziak}

\end{document}
